\section{Vorgehen und Aufbau der Arbeit}

Die Arbeit folgt einem strukturierten Vorgehen von der konzeptionellen Grundlage
über die Strategieentwicklung bis zur praktischen Anwendung.
Kapitel 2 legt die fachliche Basis mit zentralen Begriffen, Risikotypen und Governance-Strukturen.
Kapitel 3 stellt die methodischen Bausteine vor: Risikobeurteilung, -minimierung und -monitoring.
Kapitel 4 beschreibt die entwickelte RM-Extension und ihre Integration in die DASC-PM-Phasen.
Kapitel 5 überprüft die Strategie anhand zweier Fallstudien,
Kapitel 6 diskutiert die Ergebnisse und zeigt Limitationen sowie weiterführende Fragestellungen auf.
